% =====================================================================
%  interfaces.tex
%  Interfaces de usuario (v1): pantallas del prototipo RedPeppers ERP.
% =====================================================================

\section{Interfaces de usuario}
\label{sec:interfaces}

\noindent{\small\textbf{Versión 1} \textperiodcentered{} documento de trabajo sujeto a validación con el interesado.}\par\vspace{0.5em}

Las Figuras~\ref{fig:dashboard} y~\ref{fig:pos} reproducen dos pantallas
del prototipo de interfaz presentado por el equipo (ERP BackOffice v2.4).
El prototipo \textbf{no es un sistema implementado}: es una propuesta visual
para discutir la interacción y la disposición de la información con Jerry.
Por tanto, acredita intención de diseño, no funcionamiento ni cumplimiento
normativo.

La Figura~\ref{fig:dashboard} corresponde a la vista gerencial. Ordena la
jornada en cuatro bloques: el desglose del margen bruto diario —ventas
confirmadas menos la cuota prorrateada y las compras directas—, los
indicadores del día, el motor de prorrateo por insumo y proveedor, y las
compras adicionales. Al costado reúne las ventas por canal y el estado de
comprobantes ante SUNAT.

La Figura~\ref{fig:pos} corresponde al terminal de mostrador. Permite elegir
el canal de atención, filtrar por categoría y agregar productos al pedido;
el carrito muestra adicionales, cantidades y el total con IGV, y el botón
final envía el pedido a cocina. El producto sin stock aparece bloqueado, en
línea con el requisito R-ENT2-05.

\paragraph{Observación de coherencia.} El panel rotula el reparto como
«5 días» y «Prorroteo 5 días». Conforme a P-14, el prorrateo debe hacerse
sobre \emph{jornadas hábiles efectivas}, no sobre una semana fija de cinco
días; así se corrigió en los diagramas 2 y 3. El prototipo debe alinearse
con esa regla antes de considerarse definitivo. Del mismo modo, los importes
mostrados son datos de ejemplo y no cifras reales del negocio.

% La captura es una imagen de pantalla; si se cuenta con el prototipo
% navegable, conviene exportar las vistas en formato vectorial.
\begin{figure}[H]
  \centering
  \includegraphics[width=\linewidth,height=0.90\textheight,keepaspectratio]{01-dashboard-margen.png}
  \caption{Interfaz 1: panel gerencial de margen bruto y operaciones (v1).
  Presenta el desglose del margen diario, los indicadores de la jornada, el
  motor de prorrateo semanal por insumo y las compras adicionales del día.}
  \label{fig:dashboard}
\end{figure}

\begin{figure}[H]
  \centering
  \includegraphics[width=\linewidth]{02-pos-mostrador.png}
  \caption{Interfaz 2: terminal de mostrador (v1). Selección de productos
  por categoría, carrito con adicionales y cantidades, total con IGV y
  envío del pedido a cocina; el producto agotado se muestra bloqueado.}
  \label{fig:pos}
\end{figure}
